\chapter{Huber's Robust Location Estimator}
\label{sec:p15-huber-robust-location}

Huber's original location problem asks how to estimate a center when
the prototype distribution is only approximately known. Its solution
uses a convex residual criterion with a clipped derivative, and a
normal density continued by exponential tails as the least-favorable
law \cite[Secs.~1--7]{Huber1964RobustLocation}.
This chapter supplies the complete known-scale contaminated-normal
construction, its asymptotic hypotheses and its minimax meaning.
The criterion used to define an estimate and the loss used to assess
its statistical risk are different functions.

\section{Convex residual criteria}
\label{sec:p15-convex-location-estimation}

\begin{hbdefinition}{def:p15-convex-location-criterion}{A set of minimizing locations}
Let \(\rho:\mathbb R\to\mathbb R\) be continuous, convex, and
\(\rho(t)\to\infty\) as \(t\to\pm\infty\).
For an integer \(n\geq1\) and real observations \(x_1,\ldots,x_n\), put
\[
 Q_n(t)=\sum_{i=1}^n\rho(x_i-t),\qquad
 \mathcal T_n(x)=\operatorname*{arg\,min}_{t\in\mathbb R}Q_n(t).
\]
Write \(\psi_-,\psi_+\) for the left and right derivatives of \(\rho\),
and use the source's midpoint convention
\(\psi=(\psi_-+\psi_+)/2\). Thus \(\psi\) is finite and nondecreasing.
When a single estimate is needed, use the midpoint \(T_n\) of
\(\mathcal T_n\). A translation-equivariant estimator satisfies
\(T_n(x_1+a,\ldots,x_n+a)=T_n(x)+a\) for all real \(a\).
\end{hbdefinition}

\begin{hbtheorem}{thm:p15-convex-location-minimizers}{proved=annex:p15.convex.location.minimizers}{Existence, measurability and the correct derivative conditions}
The minimizing set is a nonempty compact interval. It is a singleton
if \(\rho\) is strictly convex. Its endpoints and midpoint are
Borel measurable functions of the sample, and the midpoint is
translation equivariant. A location \(t\) minimizes \(Q_n\) exactly when
\[
 \sum_i\psi_-(x_i-t)\leq0\leq\sum_i\psi_+(x_i-t).
\]
If \(\psi\) is continuous this is the equation
\(\sum_i\psi(x_i-t)=0\).
For any choice \(T\in\mathcal T_n\), put \(U_n(t)=\sum_i\psi(x_i-t)\).
Then \(U_n(t)<0\) implies \(T\leq t\), and \(T<t\) implies
\(U_n(t)\leq0\). These implications remain valid with jumps or
an interval of minimizing locations.
\end{hbtheorem}

\begin{hbtheorem}{thm:p15-convex-location-consistency}{proved=annex:p15.convex.location.consistency}{The population score and strong consistency}
Let \(X_i\) be i.i.d. real variables with law \(F\), and define
\(\lambda(t)=E_F\psi(X-t)\). If this expectation is finite at some
\(t_0\), it is well defined, possibly infinite, at every \(t\);
it is nonincreasing and is strictly positive for sufficiently negative
\(t\), strictly negative for sufficiently positive \(t\).
If a finite \(c\) satisfies
\[
 \lambda(t)>0\ (t<c),\qquad \lambda(t)<0\ (t>c),
\]
then every measurable choice \(T_n\in\mathcal T_n\) converges to \(c\)
almost surely and in probability. This is the source's Lemmas 2--3,
with the extended-expectation strong-law step made explicit.
\end{hbtheorem}

\begin{hbtheorem}{thm:p15-convex-location-normality}{proved=annex:p15.convex.location.normality}{The full convex-score normal limit}
In the i.i.d. setting of \cref{thm:p15-convex-location-consistency},
suppose \(\lambda(c)=0\), \(\lambda\) is differentiable at \(c\),
\(\lambda'(c)=-a<0\), and
\[
 B(t)=E_F\psi(X-t)^2
\]
is finite on a neighborhood of \(c\) and continuous at \(c\).
Write \(B=B(c)\). Then every measurable minimizing choice satisfies
\[
 \sqrt n(T_n-c)\Rightarrow N(0,B/a^2).
\]
If \(B=0\), this means convergence in probability to zero.
If instead the left and right derivatives of \(\lambda\) at \(c\)
are \(-a_-\), \(-a_+\), both finite and strictly negative, and \(B>0\),
the limit has the continuous distribution function
\[
 G(z)=
 \begin{cases}
 \Phi(a_-z/\sqrt B),&z<0,\\
 \Phi(a_+z/\sqrt B),&z\geq0.
 \end{cases}
\]
Here \(\Phi\) is the standard normal CDF. This is the source's
two-half-normal extension; it need not be a centered normal law.
The annex preserves its triangular-array argument with complete
Lindeberg and threshold-atom qualifications.
\end{hbtheorem}

\begin{hbtheorem}{thm:p15-nonconvex-location-qualification}{proved=annex:p15.nonconvex.location.qualification}{A conditional smooth-score theorem and a consistency failure}
Let \(X_i\) be i.i.d. real variables with law \(F\). Suppose measurable estimates satisfy
\(\sum_i\psi(X_i-T_n)=0\) and \(T_n\to0\) in probability.
Assume \(\psi\) is continuously differentiable, \(\psi'\) is uniformly
continuous on \(\mathbb R\), \(E\psi(X)^2<\infty\),
\(E|\psi'(X)|<\infty\), and \(E\psi'(X)=a>0\).
Then \(E\psi(X)=0\) and
\(\sqrt n T_n\Rightarrow N(0,E\psi(X)^2/a^2)\).
Consistency is a premise, not a consequence of smoothness alone.
For example \(\rho(t)=(t^2-1)^2\) and \(X_i=0\) almost surely give
minimizers \(\{-1,1\}\); choosing \(T_n=(-1)^n\) is inconsistent.
\end{hbtheorem}

\section{The clipped criterion and its computation}
\label{sec:p15-huber-location-computation}

\begin{hbdefinition}{def:p15-huber-location-score}{A positive clipping constant and a known scale}
For \(k>0\), define
\[
 \rho_k(x)=
 \begin{cases}x^2/2,&|x|\leq k,\\ k|x|-k^2/2,&|x|>k,\end{cases}
 \qquad
 \psi_k(x)=\max(-k,\min(k,x)).
\]
The estimate \(T_{n,k}\) is the midpoint of the zeros of
\(U(t)=\sum_i\psi_k(x_i-t)\), equivalently of the minimizing interval.
For a known scale \(\sigma>0\), use
\(\sigma T_{n,k}(x_1/\sigma,\ldots,x_n/\sigma)\).
The resulting threshold for residuals in the original units is \(k\sigma\).
Unknown scale is a separate problem.
\end{hbdefinition}

\begin{hbtheorem}{thm:p15-huber-location-computation}{proved=annex:p15.huber.location.computation}{Finite computation and its Winsorizing identity}
The score sum is continuous and nonincreasing; its zeros lie between
\(\min_i x_i\) and \(\max_i x_i\).
Sort the finite breakpoints \(x_i-k,x_i+k\).
On an interval between consecutive breakpoints let \(I\) be the
indices with \(|x_i-t|<k\), and let \(n_-,n_+\) count values below
\(t-k\) and above \(t+k\). Then
\[
 U(t)=\sum_{i\in I}x_i-|I|t+k(n_+-n_-).
\]
If \(|I|>0\), its only possible zero is
\[
 t=\{\sum_{i\in I}x_i+k(n_+-n_-)\}/|I|.
\]
Retain it if it lies in the closed interval.
If \(|I|=0\), the entire interval consists of zeros exactly when
\(n_+=n_-\). Include zero breakpoints and take the midpoint of the
resulting zero interval. This finite procedure computes the defined
estimate, including ties and flat intervals.

For any zero \(t_0\), replace every observation outside
\([t_0-k,t_0+k]\) by its nearer endpoint. The score still vanishes
at \(t_0\), and \(t_0\) is the arithmetic mean of the modified sample.
For \((0,0,0,M)\), \(M>4/3\), \(k=1\), the estimate is \(1/3\),
whereas the mean is \(M/4\) and the median is zero.
For \((-3,3)\), \(k=1\), the minimizing interval is \([-2,2]\).
\end{hbtheorem}

\section{The symmetric normal-contamination model}
\label{sec:p15-huber-normal-neighborhood}

\begin{hbdefinition}{def:p15-huber-symmetric-location-model}{The parameter and the limiting-risk order}
Let \(\phi(x)=(2\pi)^{-1/2}\exp(-x^2/2)\).
For a fixed \(0\leq\varepsilon<1\), let
\[
 \mathcal C_\varepsilon
   =\{F=(1-\varepsilon)\Phi+\varepsilon H:
                         H\text{ is any probability symmetric about }0\}.
\]
This includes discrete, continuous, singular and heavy-tailed \(H\).
Observations have law \(F(x-\xi)\), independently, and \(\xi\) is
the center of symmetry; the normal component's scale is known to be one.
For translation-equivariant measurable estimator sequences \(T=(T_n)\),
the source's limiting-risk criterion is
\[
 \mathcal R_\varepsilon(T)=
 \sup_{F\in\mathcal C_\varepsilon}
 \lim_{C\to\infty}\liminf_{n\to\infty}
             E_F\min(nT_n^2,C^2).
\]
The expectation is evaluated at center zero, using equivariance.
The supremum over \(F\) is outside the displayed limiting operations.
This criterion is not a finite-sample supremum of mean squared errors.
\end{hbdefinition}

\begin{hbtheorem}{thm:p15-huber-normal-variance}{proved=annex:p15.huber.normal.variance}{Every symmetric contaminant, including cutoff atoms}
For \(F\in\mathcal C_\varepsilon\) and fixed \(k>0\), the estimator
is strongly consistent for zero and
\[
 \sqrt n T_{n,k}\Rightarrow N(0,V_k(F)),\qquad
 V_k(F)=\frac{E_F\min(X^2,k^2)}{A_F(k)^2},
\]
where
\[
 A_F(k)=F((-k,k))+\tfrac12F(\{-k,k\})>0.
\]
The half weight at the endpoints is necessary for symmetric atoms there.
Put
\[
 A(k)=2\Phi(k)-1,\quad
 B(k)=A(k)-2k\phi(k)+2k^2(1-\Phi(k)).
\]
Then the exact worst asymptotic variance is
\[
 V^*(k,\varepsilon)=
 \frac{(1-\varepsilon)B(k)+\varepsilon k^2}
      {(1-\varepsilon)^2A(k)^2}.
\]
It is attained by every symmetric \(H\) supported on \(\{|x|>k\}\).
At known scale \(\sigma\), the asymptotic variance is multiplied by
\(\sigma^2\).
\end{hbtheorem}

\begin{hbtheorem}{thm:p15-huber-least-favorable-normal}{proved=annex:p15.huber.least.favorable.normal}{Normalization, Fisher information and the score saddlepoint}
For \(0<\varepsilon<1\), there is a unique \(k>0\) satisfying
\[
 (1-\varepsilon)^{-1}=A(k)+2\phi(k)/k.
\]
The density
\[
 f_*(x)=(1-\varepsilon)(2\pi)^{-1/2}\exp(-\rho_k(x))
\]
integrates to one and equals
\((1-\varepsilon)\phi(x)\) on \([-k,k]\), with exponential tails
outside. It belongs to \(\mathcal C_\varepsilon\), is positive and
continuously differentiable, and has location score
\(-f_*'/f_*=\psi_k\). Its Fisher information is
\[
 I_*=\int(f_*'/f_*)^2f_*=(1-\varepsilon)A(k),
 \qquad V^*(k,\varepsilon)=1/I_*.
\]
For any measurable score \(\psi\) with \(E_{F_*}\psi=0\),
\(0<E_{F_*}\psi^2<\infty\), and
\(a_\psi=-\frac{d}{dt}E_{F_*}\psi(X-t)|_{t=0}\ne0\),
\[
 \frac{E_{F_*}\psi^2}{a_\psi^2}\geq1/I_*.
\]
Equality holds exactly when \(\psi\) is a nonzero multiple of \(\psi_k\)
\(F_*\)-almost everywhere. The derivative and integration identity
used here are proved in the annex.
The density \(f_*\) uniquely minimizes Fisher information among
contamination densities \(f\) that are absolutely continuous as
functions on \(\mathbb R\) and have finite
\(I(f)=\int_{\{f>0\}}f'^2/f\).
\end{hbtheorem}

\begin{hbtheorem}{thm:p15-location-integrated-score-bound}{proved=annex:p15.location.integrated.score.bound}{The exact lower bound needed for all equivariant estimators}
Let \(f>0\) be a continuously differentiable probability density,
with \(I(f)=\int f'^2/f<\infty\), and let the observations be
independent with common location density \(f(x-\theta)\).
For any real-valued measurable
translation-equivariant estimate \(T_n\), and every \(C>0\),
\[
 E_f\min(nT_n^2,C^2)\geq
                 \frac1{I(f)+4\pi^2/C^2}.
\]
The annex proves this by integration over an explicit shrinking
location prior and a clipped estimator. No finite second moment
of \(T_n\), unbiasedness or limiting interchange is assumed.
\end{hbtheorem}

\begin{hbtheorem}{thm:p15-huber-location-minimax}{proved=annex:p15.huber.location.minimax}{The original all-translation-equivariant minimax conclusion}
For \(0<\varepsilon<1\) and the normalized \(k\) above,
\[
 \inf_{T\ {\rm translation\ equivariant}}
        \mathcal R_\varepsilon(T)
 =\mathcal R_\varepsilon(T_{\cdot,k})
 =V^*(k,\varepsilon)=1/I_*.
\]
The source's score-variance saddlepoint and the stronger all-equivariant
claim are therefore both supplied with complete proofs.
At \(\varepsilon=0\), the sample mean attains the corresponding value one.
The formal substitution \(k=0\) makes \(\rho_k\) constant and does
not define a median estimator; the median arises from the limiting
normalized criterion \(\rho_k/k\to|x|\) as \(k\downarrow0\).
\end{hbtheorem}

\section{Bias, finite risks and a general-information boundary}
\label{sec:p15-huber-location-boundaries}

\begin{hbtheorem}{thm:p15-huber-asymmetric-bias}{proved=annex:p15.huber.asymmetric.bias}{An exact contamination-bias equation}
Fix \(k>0\), \(0\leq\varepsilon<1\), and any probability \(H\) on
\(\mathbb R\). The law
\(F=(1-\varepsilon)\Phi+\varepsilon H\) has a unique population
score zero \(c(F)\), and the estimate converges to it.
For \(0\leq\varepsilon<1/2\), the exact maximum absolute displacement
is \(c_*\geq0\), determined by
\[
 m_k(c_*)=-\frac{\varepsilon k}{1-\varepsilon},\qquad
 m_k(c)=E_\Phi\psi_k(Z-c).
\]
The positive displacement is attained by contamination at any
\(M>c_*+k\). Here
\[
\begin{split}
 m_k(c)={}&\phi(c-k)-\phi(c+k)
 -c\{\Phi(c+k)-\Phi(c-k)\}\\
 &+k\{1-\Phi(c+k)-\Phi(c-k)\}.
\end{split}
\]
For fixed \(k\), as \(\varepsilon\downarrow0\),
\[
 c_*=\frac{\varepsilon k}{(1-\varepsilon)A(k)}
          +O\!\left((\varepsilon/(1-\varepsilon))^3\right).
\]
For \(1/2\leq\varepsilon<1\), the supremum of \(c(F)\) is infinite.
Without symmetry, this displacement concerns the stipulated normal
component; the component's location need not be identified from the
mixture law alone.
\end{hbtheorem}

\begin{hbtheorem}{thm:p15-huber-location-failure-cases}{proved=annex:p15.huber.location.failure.cases}{Cutoff atoms, variance inflation and the order of limits}
Let \(0<\varepsilon<1\) and
\(H=(\delta_k+\delta_{-M})/2\), \(M>k\).
Then \(E_F\psi_k=0\), but the two population-score slopes at zero are
\[
 a_-=(1-\varepsilon)A(k),\qquad
 a_+=(1-\varepsilon)A(k)+\varepsilon/2.
\]
The resulting limit is the two-half-normal law, not a single centered
normal law.

For symmetric point contamination \(H=(\delta_M+\delta_{-M})/2\),
the sample mean's asymptotic variance is \(1-\varepsilon+\varepsilon M^2\).
For \(M>k\), the clipped estimate instead has the fixed asymptotic
variance \(V^*(k,\varepsilon)\).
Nevertheless for every fixed sample size \(n\) and \(\varepsilon>0\),
the supremum of the actual finite-sample variance of any estimator
always lying in the sample's convex hull is infinite over symmetric
contamination. Thus the limiting-risk order in the minimax theorem
cannot be replaced by a finite-sample supremum.
\end{hbtheorem}

\begin{hbtheorem}{thm:p15-fisher-minimizer-nonuniqueness}{proved=annex:p15.fisher.minimizer.nonuniqueness}{Convexity alone does not give a unique information minimizer}
Let \(f(x)=2\sin^2(\pi x)\) on \(0<x<1\), zero elsewhere,
and \(g(x)=f(x-2)\). Here the vague topology is defined by integration
against every continuous compactly supported function.
The convex class of densities
\(\{(1-t)f+tg:0\leq t\leq1\}\) is compact for the vague topology.
Every density is absolutely continuous and has the same finite
Fisher information \(4\pi^2\).
Thus a general convex compact class need not have a unique Fisher
information minimizer. This does not affect the uniqueness of the
positive normal-contamination density proved above.
\end{hbtheorem}

The paper's general log-concave prototype theorem and complete abstract
information game, including the support and uniqueness corrections, are
supplied in \cref{sec:p15-huber-information-game}.
The alternative CDF neighborhood, scale estimation, all three joint
location/scale proposals, the \(M_r\) extension and the historical
numerical tables are supplied in \cref{sec:p15-huber-remaining-methods}.
The full joint law at clipping atoms and the scale comparison allowing
observation signs are supplied in \cref{sec:p15-huber-joint-endpoints}.
Later influence-function and breakdown theories are separate original
sources; they are not attributed to this 1964 result.
